\documentclass[11pt,letterpaper]{article}
\usepackage[margin=1in]{geometry}
\usepackage{hyperref}
\usepackage{enumitem}
\usepackage{titlesec}
\usepackage{parskip}
\usepackage{booktabs}
\usepackage{array}
\usepackage{xcolor}

\hypersetup{colorlinks=true, urlcolor=blue, linkcolor=blue}

\title{\textbf{Research Paper Summaries}\\
\large Layers, Security Skills, and Technical Expertise}
\author{Jaehyuk Lee}
\date{}

\begin{document}
\maketitle
\tableofcontents
\newpage

% ============================================================
\section*{Overview: Skills Across All Papers}
\addcontentsline{toc}{section}{Overview: Skills Across All Papers}
% ============================================================

The nine papers collectively span four major \textbf{system layers} and two broad \textbf{security skill domains}:

\subsection*{Layers Touched}
\begin{itemize}
    \item \textbf{Hardware / Microarchitecture:} CPU cache internals, DRAM (Rowhammer), dirty-bit logic, PMU counters, cache replacement policy (Tree-PLRU), Intel MEE, SGX EPC.
    \item \textbf{Firmware / TEE Runtime:} ARM TF-RMM, TrustZone Monitor mode, SGX enclave lifecycle instructions (ECREATE, EADD, EENTER, EEXIT), SMMU.
    \item \textbf{Hypervisor / OS Kernel:} KVM, Stage-2 page tables, cgroups, Linux process management, QEMU binary translation, perf/PMU kernel interface.
    \item \textbf{Application / Protocol:} SGX enclave user-space library (\texttt{sgxlib}), Solidity smart contracts, DeFi protocol logic, fuzzer runtimes (AFL, libFuzzer).
\end{itemize}

\subsection*{Security Skills Learned}
\begin{itemize}
    \item \textbf{TEE design and exploitation:} Building, attacking, and defending Trusted Execution Environments (SGX, TrustZone, Arm CCA).
    \item \textbf{Microarchitectural side-channel attacks:} Cache timing (Prime+Probe, Prime+Retouch), dirty-bit exploitation, PMU-based leakage.
    \item \textbf{Hardware fault attacks:} Rowhammer against SGX EPC, MEE integrity fault as a denial-of-service primitive.
    \item \textbf{Control-flow attacks:} Return-Oriented Programming (ROP) against TrustZone Secure World.
    \item \textbf{DMA isolation and peripheral security:} Extending CCA's isolation guarantees to SMMU and DMA-capable peripherals.
    \item \textbf{Automated vulnerability discovery:} Coverage-guided fuzzing, dynamic fuzzer composition, runtime scheduling for exploit discovery.
    \item \textbf{Smart contract / DeFi security:} Vulnerability taxonomy, longitudinal exploit analysis, on-chain forensics.
\end{itemize}

\newpage

% ============================================================
\section{PORTAL: Fast and Secure Device Access with Arm CCA}
% ============================================================
\textbf{Venue:} IEEE S\&P 2025 \hfill
\href{https://www.computer.org/csdl/proceedings-article/sp/2025/223600a013/21B7Q5c3ZHq}{[Paper Link]}

\subsection*{Problem}
Modern mobile SoCs integrate peripherals (GPU, NPU, DSP) that use DMA to access memory.
When confidential VMs run under Arm CCA (Confidential Compute Architecture), these peripherals
can break the Realm isolation boundary because the IOMMU/SMMU is not CCA-aware by design.

\subsection*{What Was Built}
PORTAL extends Arm CCA's Realm Management Monitor (TF-RMM) to mediate peripheral device access.
It introduces a new interface between the RMM and the normal-world hypervisor (KVM) so that
Realm VMs can safely share DMA with peripherals without compromising confidentiality or integrity.

\subsection*{Layers Touched}
\begin{itemize}
    \item \textbf{Firmware:} TF-RMM — extended to track Realm memory regions and enforce SMMU mappings.
    \item \textbf{Hypervisor:} KVM — new RMI calls to request/release device access on behalf of Realm VMs.
    \item \textbf{Driver:} SMMU driver — modified to enforce Realm-owned PA ranges.
    \item \textbf{Guest VM:} Realm guest driver — userspace API to request DMA from peripherals.
\end{itemize}

\subsection*{Security Skills}
\begin{itemize}
    \item Arm CCA architecture and Realm Management Interface (RMI)
    \item Stage-2 page table design for confidential VMs
    \item SMMU/IOMMU internals and DMA isolation enforcement
    \item TF-RMM firmware development
    \item Linux KVM hypervisor modification
\end{itemize}

\newpage

% ============================================================
\section{A Longitudinal Study of Vulnerabilities in the DeFi Ecosystem}
% ============================================================
\textbf{Venue:} USENIX Security 2024

\subsection*{Problem}
DeFi protocols have suffered billions of dollars in losses from smart contract exploits.
No prior work systematically studied how vulnerability classes evolve over time and
whether the ecosystem learns from past incidents.

\subsection*{What Was Built}
A dataset and analysis pipeline covering hundreds of DeFi incidents over multiple years,
categorizing vulnerability classes (reentrancy, price oracle manipulation, flash loan attacks,
access control flaws, etc.) and tracking recurrence trends.
Qualitative coding was performed using Dedoose.

\subsection*{Layers Touched}
\begin{itemize}
    \item \textbf{Smart contract layer:} Solidity code auditing and vulnerability classification.
    \item \textbf{Protocol logic:} AMMs, lending protocols, flash loan mechanics.
    \item \textbf{On-chain data:} Blockchain transaction analysis and incident reconstruction.
\end{itemize}

\subsection*{Security Skills}
\begin{itemize}
    \item Smart contract vulnerability taxonomy and auditing
    \item Blockchain data analysis and on-chain forensics
    \item DeFi protocol mechanics (AMMs, flash loans, oracles)
    \item Qualitative research methodology (Dedoose coding)
    \item Longitudinal security measurement methodology
\end{itemize}

\newpage

% ============================================================
\section{SENSE: Side-Channel Evaluation and Noise-addition System for Enclaves}
% ============================================================
\textbf{Venue:} NDSS 2024

\subsection*{Problem}
SGX enclaves are vulnerable to microarchitectural side-channel attacks (cache timing,
branch predictor leakage, PMU-based observation). Defenses are ad hoc and lack a
systematic evaluation framework.

\subsection*{What Was Built}
SENSE provides: (1) a systematic framework for evaluating microarchitectural side-channel
leakage inside and across SGX enclave boundaries, and (2) a noise-injection mechanism
to mitigate observable leakage. It profiles hardware performance counters inside and
outside the enclave using the kernel perf subsystem.

\subsection*{Layers Touched}
\begin{itemize}
    \item \textbf{CPU microarchitecture:} Cache behavior, branch predictor state, PMU event counting.
    \item \textbf{SGX enclave runtime:} Enclave entry/exit (AEX), SSA-based state observation.
    \item \textbf{OS kernel:} Perf subsystem and PMU access from ring-0.
    \item \textbf{Compiler/toolchain:} Instrumentation of enclave binaries.
\end{itemize}

\subsection*{Security Skills}
\begin{itemize}
    \item Intel SGX internals (AEX, SSA, EPCM)
    \item Microarchitectural side-channel attack methodology
    \item PMU/perf event profiling at kernel level
    \item GEM5 simulation for microarchitectural analysis
    \item Side-channel leakage quantification and noise mitigation
\end{itemize}

\newpage

% ============================================================
\section{autofz: Automated Fuzzer Composition at Runtime}
% ============================================================
\textbf{Venue:} USENIX Security 2023 \hfill
\href{https://github.com/sslab-gatech/autofz}{[GitHub]}

\subsection*{Problem}
Dozens of fuzzers exist, each optimized for different target characteristics.
Manual fuzzer selection is suboptimal, and static combinations waste resources.
No prior system dynamically allocates fuzzer compute resources based on live feedback.

\subsection*{What Was Built}
autofz monitors coverage feedback (AFL bitmap) from multiple fuzzers running in parallel
and applies an AIMD-inspired (Additive Increase, Multiplicative Decrease) scheduler to
dynamically reallocate CPU time toward whichever fuzzers are currently finding new coverage.

\subsection*{Layers Touched}
\begin{itemize}
    \item \textbf{OS scheduling:} cgroups for CPU quota enforcement, process pinning.
    \item \textbf{Fuzzer internals:} AFL bitmap coverage feedback, libFuzzer corpus sharing.
    \item \textbf{Coverage instrumentation:} Compile-time and binary-level coverage tracking.
\end{itemize}

\subsection*{Security Skills}
\begin{itemize}
    \item Coverage-guided fuzzing (AFL, libFuzzer, honggfuzz, etc.)
    \item Linux cgroups and process resource management
    \item AIMD scheduling algorithm design
    \item Automated vulnerability discovery methodology
    \item Runtime profiling and feedback-driven optimization
\end{itemize}

\newpage

% ============================================================
\section{Hacking in Darkness: ROP against TrustZone Secure World}
% ============================================================
\textbf{Venue:} USENIX Security 2017

\subsection*{Problem}
ARM TrustZone's Secure World is assumed to be a strong isolation boundary.
This paper demonstrates that ROP (Return-Oriented Programming) attacks can be
mounted \emph{against} the Secure World from the Normal World, breaking that assumption.

\subsection*{What Was Built}
A full ROP attack chain targeting a real TrustZone implementation.
Gadgets were identified in Secure World firmware binaries, chained via SVC/SMC calls
from Normal World Linux, and used to achieve arbitrary code execution inside Secure World.

\subsection*{Layers Touched}
\begin{itemize}
    \item \textbf{ARM TrustZone:} Monitor mode, Secure EL1 firmware, SVC/SMC dispatch.
    \item \textbf{Normal World OS:} Linux userspace and kernel as the attacker's vantage point.
    \item \textbf{Binary analysis:} Gadget discovery in Secure World binaries.
\end{itemize}

\subsection*{Security Skills}
\begin{itemize}
    \item Return-Oriented Programming (ROP) exploit construction
    \item ARM TrustZone architecture (Monitor mode, world switching)
    \item Binary analysis and gadget finding
    \item SVC/SMC call interface exploitation
    \item QEMU-based system emulation for attack development
\end{itemize}

\newpage

% ============================================================
\section{SGX-Bomb: Locking Down the Processor via Rowhammer}
% ============================================================
\textbf{Venue:} SysTEX 2017

\subsection*{Problem}
Intel SGX's Memory Encryption Engine (MEE) detects EPC bit flips to protect enclave integrity.
This defense can be weaponized: an attacker who induces Rowhammer bit flips in EPC causes
the MEE to trigger a system-wide processor lockdown, creating a denial-of-service.

\subsection*{What Was Built}
SGX-Bomb demonstrates that Rowhammer against EPC pages causes the MEE integrity check to fail,
resulting in a machine check exception that locks the processor and requires a physical reboot.
This is a novel \emph{liveness attack} using SGX's own security mechanism against the system.

\subsection*{Layers Touched}
\begin{itemize}
    \item \textbf{DRAM hardware:} Rowhammer bit-flip induction via cache eviction hammering.
    \item \textbf{Intel SGX:} EPC memory layout, MEE integrity verification, EPCM state.
    \item \textbf{CPU fault handling:} Machine check exceptions triggered by MEE failure.
\end{itemize}

\subsection*{Security Skills}
\begin{itemize}
    \item Rowhammer attack methodology (cache flush, hammering patterns)
    \item Intel SGX internals (EPC, MEE, EPCM)
    \item Hardware fault injection and exploitation
    \item Denial-of-service via microarchitectural fault escalation
    \item Memory subsystem security analysis
\end{itemize}

\newpage

% ============================================================
\section{PrivateZone: Providing a Private Execution Environment using ARM TrustZone}
% ============================================================
\textbf{Venue:} IEEE TDSC 2016

\subsection*{Problem}
Sensitive mobile applications lack hardware-backed isolation equivalent to Intel SGX.
ARM TrustZone exists but is locked to OEM/vendor code and unavailable to third-party apps.

\subsection*{What Was Built}
PrivateZone opens TrustZone's Secure World to third-party trusted applications by using
Stage-2 paging and Monitor mode to isolate Secure World apps from each other and from
the Normal World — providing an SGX-like enclave model on ARM mobile platforms.

\subsection*{Layers Touched}
\begin{itemize}
    \item \textbf{ARM TrustZone:} Monitor mode, Secure EL1 OS design, world-switch handler.
    \item \textbf{Stage-2 paging:} Memory isolation between multiple Secure World tenants.
    \item \textbf{Normal World OS:} Linux kernel interface for launching trusted apps.
    \item \textbf{Emulation:} QEMU modifications to support PrivateZone semantics.
\end{itemize}

\subsection*{Security Skills}
\begin{itemize}
    \item ARM TrustZone design and Secure World OS implementation
    \item Stage-2 (IPA$\to$PA) page table construction for tenant isolation
    \item Monitor mode entry/exit and world-switch security
    \item Trusted application lifecycle management
    \item QEMU system modification for TEE research
\end{itemize}

\newpage

% ============================================================
\section{OpenSGX: An Open Platform for SGX Research}
% ============================================================
\textbf{Venue:} NDSS 2015

\subsection*{Problem}
Intel SGX was announced before real hardware was available, leaving researchers
with no way to build, run, or study SGX enclaves. A faithful, open emulation platform was needed.

\subsection*{What Was Built}
OpenSGX is a full software emulation of Intel SGX built on QEMU.
It emulates SGX instructions (ECREATE, EADD, EENTER, EEXIT, ERESUME),
EPCM, SECS, TCS, and SSA data structures.
It also ships \texttt{sgxlib}, a user-space library for writing enclave programs,
and adds GDB support for in-enclave debugging.

\subsection*{Layers Touched}
\begin{itemize}
    \item \textbf{CPU emulation:} QEMU binary translation extended with SGX opcode handling.
    \item \textbf{SGX ISA:} All enclave lifecycle instructions and memory management structures.
    \item \textbf{User-space runtime:} \texttt{sgxlib} enclave programming library.
    \item \textbf{Debug tooling:} GDB extension for single-stepping inside emulated enclaves.
\end{itemize}

\subsection*{Security Skills}
\begin{itemize}
    \item Intel SGX ISA and enclave lifecycle (ECREATE$\to$EENTER$\to$EEXIT)
    \item QEMU internals and binary translation extension
    \item Enclave memory management (EPC, EPCM, SECS, TCS, SSA)
    \item Systems programming for TEE runtime design
    \item Debugger integration for security research tooling
\end{itemize}

\newpage

% ============================================================
\section{PRIME+RETOUCH: Exploiting Cache Dirty Bits for Stealthy Cache Attacks}
% ============================================================
\textbf{Venue:} arXiv

\subsection*{Problem}
Classic Prime+Probe cache side-channel attacks are noisy because cache eviction is
detectable by victim or defense monitors.
The write-back cache dirty-bit mechanism is an under-explored side channel that leaks
cache-line usage without requiring the noisy eviction step.

\subsection*{What Was Built}
PRIME+RETOUCH exploits write-back dirty bits to perform stealthy cache side-channel attacks.
By ``retouching'' (re-dirtying) cache lines, the attacker avoids eviction-based detection
used by defenses that monitor cache evictions.
The attack is evaluated on Intel and ARM M1 platforms using hardware PMUs
and includes a detailed analysis of the Tree-PLRU cache replacement policy.

\subsection*{Layers Touched}
\begin{itemize}
    \item \textbf{CPU cache microarchitecture:} L1/LLC write-back dirty-bit logic, cache-line state.
    \item \textbf{Cache replacement policy:} Tree-PLRU analysis to control eviction behavior.
    \item \textbf{PMU:} Hardware performance counters on Intel (TSX) and ARM M1.
\end{itemize}

\subsection*{Security Skills}
\begin{itemize}
    \item Cache side-channel attack construction (Prime+Probe variant)
    \item Write-back dirty-bit microarchitectural exploitation
    \item Tree-PLRU replacement policy reverse engineering
    \item Intel TSX abuse for cache state observation
    \item ARM M1 PMU profiling and microarchitectural timing analysis
    \item Stealth side-channel design (evading eviction-based defenses)
\end{itemize}

\newpage

% ============================================================
\section*{Consolidated Skill Matrix}
\addcontentsline{toc}{section}{Consolidated Skill Matrix}
% ============================================================

\begin{center}
\renewcommand{\arraystretch}{1.4}
\begin{tabular}{>{\bfseries}p{4cm} p{10cm}}
\toprule
Skill Area & Papers / Evidence \\
\midrule
TEE design \& impl. & OpenSGX, PrivateZone, PORTAL \\
TEE exploitation & Hacking in Darkness (ROP), SGX-Bomb (fault), SENSE (side-channel) \\
Microarch. side channels & SENSE, PRIME+RETOUCH \\
Hardware fault attacks & SGX-Bomb (Rowhammer + MEE) \\
DMA \& peripheral isolation & PORTAL (SMMU, CCA) \\
Hypervisor / Stage-2 paging & PORTAL (KVM), PrivateZone (Stage-2) \\
Firmware development & PORTAL (TF-RMM), PrivateZone (Monitor mode) \\
Fuzzing \& vuln. discovery & autofz \\
Smart contract security & DeFi longitudinal study \\
Emulation / tooling & OpenSGX (QEMU + GDB) \\
Control-flow exploitation & Hacking in Darkness (ROP gadget chaining) \\
\bottomrule
\end{tabular}
\end{center}

\end{document}
